\chapter{Maximum Likelihood Estimation: The Bradley--Terry Model}

To find and understand the underlined mechanics of the model in the choice set, we will observe some data where we need to  assume a probabilistic model, and for now we are using Bradley terry model.
Actually, the goal of doing so is to find that unknown parameter i.e. the compressed representation of the observed data.
\section{Assumption followed for Estimation}
The Assumptions that one typically works with when we are dealing with the data about the observations are following :
\[
\text{Independence:}\qquad P(X_i \mid X_j) = P(X_i), \qquad \forall\, i \neq j.
\]

\[
\text{Identically distributed:}\qquad P(X_i = 1) = P(X_j = 1) = p,
\qquad \forall\, i,j.
\]
\section{Likelihood Function}
In order to get a principled way to get estimates from the data, we need a likelihood function.
Likelihood function will give us the value of the parameter where the chance of seeing the data is highest.
As per Fisher's principle of maximum likelihood \citep{fisher1922mathematical}, given $n$
observations $x_1, x_2, \dots, x_n$ drawn from a distribution parameterised by $p$, the
likelihood function is defined as the joint probability of observing the data, viewed as a
function of the parameter:
\begin{align}
    L(p ; \{x_1, x_2, \dots, x_n\}) &= P(x_1, x_2, \dots, x_n ; p) \label{eq:likelihood-def} \\
    &= P(x_1 ; p) \cdot P(x_2 ; p) \cdots P(x_n ; p) \qquad \text{(by independence)} \label{eq:likelihood-product} \\
    &= \prod_{i=1}^{n} p^{x_i} (1-p)^{1-x_i} \label{eq:likelihood-bernoulli}
\end{align}
where the last line follows from taking each $x_i$ to be a Bernoulli trial with success
probability $p$, so that $P(x_i ; p) = p^{x_i}(1-p)^{1-x_i}$. The maximum likelihood estimate
$\hat{p}$ is then the value of $p$ that maximises Equation~\eqref{eq:likelihood-bernoulli}.
% \section*{Setting up the likelihood}

% Let us consider that we observe $N$ independent pairwise comparisons, where in the
% $n^{\text{th}}$ comparison item $i_n$ is chosen over item $j_n$. Under
% the multinomial logit model with $|C| = 2$, then, the  probability of the
% observed winner is

% \begin{equation*}
%     P(i_n \succ j_n)
%     = \frac{\exp(v_{i_n})}{\exp(v_{i_n}) + \exp(v_{j_n})}
%     = \sigma(v_{i_n} - v_{j_n})
% \end{equation*}

% where $\sigma(x) = 1 / (1 + e^{-x})$ is the logistic function. So the
% likelihood of the full data is just the product over all $N$ comparisons. 

% \begin{equation*}
%     \mathcal{L}(\mathbf{v})
%     = \prod_{n=1}^{N} \sigma(v_{i_n} - v_{j_n}).
% \end{equation*}

% Taking log for the numerical convenience,

% \begin{equation*}
%     \ell(\mathbf{v})
%     = \log \mathcal{L}(\mathbf{v})
%     = \sum_{n=1}^{N} \log \sigma(v_{i_n} - v_{j_n}).
% \end{equation*}

% \section*{Gradient of the log-likelihood}

% Now, to find the MLE we would want to maximize $\ell(\mathbf{v})$, so let
% us take the partial derivative with respect to $v_k$.

% For a single comparison $(i, j)$,

% \begin{align*}
%     \frac{\partial}{\partial v_i} \log \sigma(v_i - v_j)
%         &= 1 - \sigma(v_i - v_j) \\
%     \frac{\partial}{\partial v_j} \log \sigma(v_i - v_j)
%         &= -\bigl(1 - \sigma(v_i - v_j)\bigr)
% \end{align*}

% and summing over all $N$ comparisons,

% \begin{equation*}
%     \frac{\partial \ell}{\partial v_k}
%     = \underbrace{W_k}_{\text{wins of } k}
%     \;-\;
%     \underbrace{\sum_{n \,:\, k \in \{i_n, j_n\}} P(k \text{ beats opponent})}
%         _{\text{expected wins of } k}
% \end{equation*}

% or, in short form,

% \begin{equation*}
%     \frac{\partial \ell}{\partial v_k} = W_k - E_k.
% \end{equation*}

% This gives a very clean interpretation $\rightarrow$ at the MLE, the
% observed wins equal the expected wins for every item. That is called
% the moment matching condition.

% \section*{Gradient descent update}

% Since $\ell$ is strictly concave in $\mathbf{v}$, gradient ascent will
% converge to the unique global maximum. So the update rule becomes

% \begin{equation*}
%     v_k^{(t+1)} = v_k^{(t)} + \eta \, \bigl(W_k - E_k^{(t)}\bigr)
% \end{equation*}

% where $\eta$ is the learning rate. And after every step we also enforce
% the identifiability constraint $\sum_k v_k = 0$ by subtracting the mean
% so that the parameters are uniquely defined (same shift trick as in the
% MNL section above).

% \section*{Numerical check on the same commute example}

% Using the weights from Section 2, that is
% $w = (0.40,\, 0.25,\, 0.15,\, 0.12,\, 0.08)$ for
% $(C, T, Bi, Bu, W)$, the corresponding centered true parameters are

% \begin{equation*}
%     \mathbf{v}_{\text{true}}
%     = [\,0.853,\ 0.383,\ -0.128,\ -0.351,\ -0.757\,].
% \end{equation*}

% We simulate $N = 2000$ pairwise comparisons: for each comparison, a pair
% $(i, j)$ is drawn uniformly at random from the $5$ items, and the winner
% is drawn from Bernoulli$\bigl(\sigma(v_i - v_j)\bigr)$.

% Fitting the model in two independent ways:

% \begin{itemize}
%     \item {I-LSR} via \texttt{choix.ilsr\_pairwise} $\rightarrow$
%     Maystre's spectral method (Chapter 3 of his thesis).
%     \item {Gradient descent from scratch} $\rightarrow$ the update
%     rule above, run for $1000$ iterations with $\eta = 0.01$ on the
%     normalized gradient.
% \end{itemize}

% gives the estimates

% \begin{align*}
%     \mathbf{v}_{\text{I-LSR}} &= [\,0.827,\ 0.461,\ -0.122,\ -0.418,\ -0.749\,] \\
%     \mathbf{v}_{\text{GD}}    &\approx \mathbf{v}_{\text{I-LSR}}
% \end{align*}

% with mean absolute error $\approx 0.04$ against the true parameters.
% So, both the methods agree with each other and with the truth up to
% sampling noise. This is expected because the Bradley-Terry log-likelihood
% is strictly concave $\rightarrow$ any correct optimizer must land at the
% same unique MLE.

% Hence proved.